% GENERATED FILE. Edit canonical lessons, then run npm run generate:books.
% canonical-source-hash: fnv1a64:0bbadba510b91243
% canonical-lessons: RU-C210-ozhog, RU-C210-allergiya, RU-C210-stomatolog, RU-C210-plastyr, RU-C210-bint

\chapter{Burn, Allergy, Dentist, Sticking plaster, Bandage}
\label{ch:ru-burn-allergy-dentist-sticking-plaster-bandage}
\hlchaptermodality{210}

\begin{chapteropening}
\textbf{By the end of this chapter:} \emph{I can say a burn, an allergy, a dentist, a sticking plaster and a bandage.}

Complete the last lesson of chapter 210: I can say a burn, an allergy, a dentist, a sticking plaster and a bandage.
\end{chapteropening}

\section[\texorpdfstring{\ru{ожог} (ozhóg)}{ожог (ozhóg)}]{\ru{ожог} (ozhóg) — a burn}

\label{lesson:RU-C210-ozhog}



\begin{quote}
Before the new one: say the Russian for to joke, then the Russian for to get married.
\end{quote}

\subsection*{You'll want to know: \ru{ожог}}
\textbf{\ru{ожог}} (\emph{ozhóg}) — \textquotedblleft{}a burn\textquotedblright{}.

\begin{grammarlens}[title={What you've built}]
One more word for this chapter.
\end{grammarlens}

\subsection*{Your turn}
\begin{itemize}
\raggedright
  \item \emph{Say it:} \emph{ozhóg}
  \item \emph{Say it:} \emph{ozhóg}, once more
  \item \emph{Say it:} say \emph{ozhóg}
  \item \emph{Recall:} say the Russian for to joke, then the Russian for to get married, then say \emph{ozhóg} again
\end{itemize}

\begin{tcolorbox}[breakable,colback=teal!4,colframe=teal!35!black,title={Before you move on}]
What does \ru{ожог} mean? (\textquotedblleft{}A burn\textquotedblright{}.) Say it once more.
\end{tcolorbox}
\section[\texorpdfstring{\ru{аллергия} (allergíya)}{аллергия (allergíya)}]{\ru{аллергия} (allergíya) — an allergy}

\label{lesson:RU-C210-allergiya}



\begin{quote}
Before the new one: say the Russian for to get married, then the Russian for a burn.
\end{quote}

\subsection*{You'll want to know: \ru{аллергия}}
\textbf{\ru{аллергия}} (\emph{allergíya}) — \textquotedblleft{}an allergy\textquotedblright{}.

\begin{grammarlens}[title={What you've built}]
One more word for this chapter.
\end{grammarlens}

\subsection*{Your turn}
\begin{itemize}
\raggedright
  \item \emph{Say it:} \emph{allergíya}
  \item \emph{Say it:} \emph{allergíya}, once more
  \item \emph{Say it:} say \emph{allergíya}
  \item \emph{Recall:} say the Russian for to get married, then the Russian for a burn, then say \emph{allergíya} again
\end{itemize}

\begin{tcolorbox}[breakable,colback=teal!4,colframe=teal!35!black,title={Before you move on}]
What does \ru{аллергия} mean? (\textquotedblleft{}An allergy\textquotedblright{}.) Say it once more.
\end{tcolorbox}
\section[\texorpdfstring{\ru{стоматолог} (stomatólog)}{стоматолог (stomatólog)}]{\ru{стоматолог} (stomatólog) — a dentist}

\label{lesson:RU-C210-stomatolog}



\begin{quote}
Before the new one: say the Russian for a burn, then the Russian for an allergy.
\end{quote}

\subsection*{You'll want to know: \ru{стоматолог}}
\textbf{\ru{стоматолог}} (\emph{stomatólog}) — \textquotedblleft{}a dentist\textquotedblright{}.

\begin{grammarlens}[title={What you've built}]
One more word for this chapter.
\end{grammarlens}

\subsection*{Your turn}
\begin{itemize}
\raggedright
  \item \emph{Say it:} \emph{stomatólog}
  \item \emph{Say it:} \emph{stomatólog}, once more
  \item \emph{Say it:} say \emph{stomatólog}
  \item \emph{Recall:} say the Russian for a burn, then the Russian for an allergy, then say \emph{stomatólog} again
\end{itemize}

\begin{tcolorbox}[breakable,colback=teal!4,colframe=teal!35!black,title={Before you move on}]
What does \ru{стоматолог} mean? (\textquotedblleft{}A dentist\textquotedblright{}.) Say it once more.
\end{tcolorbox}
\section[\texorpdfstring{\ru{пластырь} (plástyr)}{пластырь (plástyr)}]{\ru{пластырь} (plástyr) — a sticking plaster}

\label{lesson:RU-C210-plastyr}



\begin{quote}
Before the new one: say the Russian for an allergy, then the Russian for a dentist.
\end{quote}

\subsection*{You'll want to know: \ru{пластырь}}
\textbf{\ru{пластырь}} (\emph{plástyr}) — \textquotedblleft{}a sticking plaster\textquotedblright{}.

\begin{grammarlens}[title={What you've built}]
One more word for this chapter.
\end{grammarlens}

\subsection*{Your turn}
\begin{itemize}
\raggedright
  \item \emph{Say it:} \emph{plástyr}
  \item \emph{Say it:} \emph{plástyr}, once more
  \item \emph{Say it:} say \emph{plástyr}
  \item \emph{Recall:} say the Russian for an allergy, then the Russian for a dentist, then say \emph{plástyr} again
\end{itemize}

\begin{tcolorbox}[breakable,colback=teal!4,colframe=teal!35!black,title={Before you move on}]
What does \ru{пластырь} mean? (\textquotedblleft{}A sticking plaster\textquotedblright{}.) Say it once more.
\end{tcolorbox}
\section[\texorpdfstring{\ru{бинт} (bint)}{бинт (bint)}]{\ru{бинт} (bint) — a bandage}

\label{lesson:RU-C210-bint}



\begin{quote}
Before the new one: say the Russian for a dentist, then the Russian for a sticking plaster.
\end{quote}

\subsection*{You'll want to know: \ru{бинт}}
\textbf{\ru{бинт}} (\emph{bint}) — \textquotedblleft{}a bandage\textquotedblright{}.

\begin{grammarlens}[title={What you've built}]
One more word for this chapter.
\end{grammarlens}

\subsection*{Your turn}
\begin{itemize}
\raggedright
  \item \emph{Say it:} \emph{bint}
  \item \emph{Say it:} \emph{bint}, once more
  \item \emph{Say it:} say \emph{bint}
  \item \emph{Recall:} say the Russian for a dentist, then the Russian for a sticking plaster, then say \emph{bint} again
\end{itemize}

\begin{tcolorbox}[breakable,colback=teal!4,colframe=teal!35!black,title={Before you move on}]
What does \ru{бинт} mean? (\textquotedblleft{}A bandage\textquotedblright{}.) Say it once more.
\end{tcolorbox}
